\documentclass[11pt]{article}
\usepackage[a4paper,margin=27mm]{geometry}
\usepackage[T1]{fontenc}
\usepackage{lmodern,microtype,amsmath,amssymb,amsthm,mathtools,enumitem,booktabs}
\usepackage[hidelinks,hypertexnames=false]{hyperref}
\newtheorem{theorem}{Theorem}
\newtheorem{lemma}[theorem]{Lemma}
\newtheorem{proposition}[theorem]{Proposition}
\newtheorem{corollary}[theorem]{Corollary}
\theoremstyle{remark}
\newtheorem{remark}[theorem]{Remark}
\newcommand{\R}{\mathbb R}
\newcommand{\Dscr}{\mathcal D}
\newcommand{\Escr}{\mathcal E}
\newcommand{\one}{\mathbf 1}
\newcommand{\supp}{\operatorname{supp}}
\newcommand{\sd}{\mathbin{\triangle}}
\newcommand{\lpr}[1]{\underline{P_{#1}}}
\newcommand{\lowprob}[1]{\underline{p_{#1}}}
\newcommand{\upprob}[1]{\overline{p_{#1}}}
\title{Independent Rankings of Binary Advice\\[3pt]\large Version 3: two-world coverage and a higher-dimensional obstruction}
\author{}
\date{9 October 2026}
\begin{document}
\maketitle

\noindent
We rank good advice above silence, silence above bad advice, and bad advice
above contradictory advice. We also require that a comparison between two
models depend only on the binary problems on which they differ.
The ordinal benchmark recipe of v2 fails this independence requirement.
A different construction, using penalties for individual binary problems,
satisfies independence and gives full coverage of coherent sets of desirable
gambles when there are two worlds. It includes boundary and vacuous targets.
For three or more worlds, no additive construction of this kind can give
full coverage: under the remaining axioms, the vacuous model must be strictly
dominated. Full coverage by more general independent ordinal rankings
remains open.

The proof has three steps. Specify penalties for the four kinds of advice;
show that their probability-weighted minima describe coherent models; then
use those minima to obtain competitiveness, closed comparison sets, and
undominated targets. The two-world construction makes every step explicit.
It uses neither benchmark scores nor a bonus depending on the whole model.
The final section explains why this method cannot give full coverage in
higher dimensions; its proof is kept in Appendix~\ref{app:obstruction}.

\section{Binary advice and the requirements}

Let $\Omega=\{\omega_1,\ldots,\omega_n\}$, $n\ge2$, and
$\Delta=\{p\in\R^n:p\ge0,\ \sum_i p_i=1\}$.
Write $e_i$ for the $i$th unit vector, $\one=(1,\ldots,1)$, and
$\supp p=\{i:p_i>0\}$. A gamble $g$ pays $g_i$ at $\omega_i$.
A model is a Lebesgue-measurable acceptance set $\Dscr\subseteq\R^n$.
Fix standard Gaussian measure $\mu$; it has the same null sets as Lebesgue
measure. Write $\Dscr\approx\Escr$ when $\mu(\Dscr\sd\Escr)=0$.
Every comparison in this note ignores changes on null sets.

A pair $\{g,-g\}$ represents a binary decision problem: if $g=f-h$, accepting
$g$ recommends $f$ over $h$, while accepting $-g$ recommends $h$ over $f$.
Neither acceptance gives no recommendation; both give contradictory
recommendations. To count each pair once, use
\[
 H=\{g:g_1>0\}.
\]
Pairs in $\{g_1=0\}$, and worlds with $g_i=0$, can be ignored because the
coordinate hyperplanes are null. The advice of $\Dscr$ at $g\in H$ is
\[
 a_{\Dscr}(g)=
 \begin{cases}
 +,&g\in\Dscr,\ -g\notin\Dscr,\\
 0,&g\notin\Dscr,\ -g\notin\Dscr,\\
 -,&g\notin\Dscr,\ -g\in\Dscr,\\
 \times,&g\in\Dscr,\ -g\in\Dscr.
 \end{cases}
\]
At $\omega_i$, $+$ is good and $-$ is bad when $g_i>0$; their roles reverse
when $g_i<0$. The background order, with better advice first, is
\begin{equation}
 \boxed{\text{good}\;\prec\;\text{no advice}\;\prec\;\text{bad}\;\prec\;\text{contradictory}.}
 \label{eq:advice-order}
\end{equation}
Write $\Dscr\unlhd_i\Escr$ when $\Dscr$ gives advice at least as good for almost
every $g\in H$. Its strict part requires strictly better advice on a
positive-measure set of pairs. The ideal acceptance set is
$D_i=\{g:g_i>0\}$.

This replaces the earlier error-inclusion order. The two orders agree on
comparisons between models that never give contradictory advice, but differ
on comparisons involving contradictions. In particular, removing the sole
correct recommendation from a contradictory pair now improves the advice,
whereas it increased the earlier error set.

A ranking consists of a preorder $\preceq_i$ at each world; smaller means
better. Its strict part is $\prec_i$. A model dominates another if it is at
least as good at every world and strictly better at some world; it strictly
dominates if it is strictly better at every world. A model is undominated
if none dominates it. It is competitive on $S\ne\varnothing$ if, against
every model, it is at least as good at some world in $S$.

We impose two requirements on comparisons:
\begin{enumerate}[leftmargin=*,itemsep=5pt]
\item[\textbf{(G)}] \textbf{Guidance improvement.} Both weak and strict
background comparisons are respected:
$\Dscr\unlhd_i\Escr\Rightarrow\Dscr\preceq_i\Escr$, and strict background
improvement implies $\Dscr\prec_i\Escr$.
\item[\textbf{(IB)}] \textbf{Independence across binary problems.} For every
measurable $A=-A$ and all measurable $F,G\subseteq A^c$,
\begin{equation}
 \begin{aligned}
 (\Dscr\cap A)\cup F&\preceq_i(\Escr\cap A)\cup F\\
 &\Longleftrightarrow\\[-6pt]
 (\Dscr\cap A)\cup G&\preceq_i(\Escr\cap A)\cup G.
 \end{aligned}\label{eq:independence}
\end{equation}
The symmetry of $A$ keeps the two orientations of each binary problem
together. Replacing shared advice outside $A$ cannot change the comparison.
\end{enumerate}

A coherent set of desirable gambles excludes $0$, contains every nonzero
$g\ge0$, and is closed under addition and positive scaling.
A credal set is a nonempty closed convex $M\subseteq\Delta$. Write
\[
 \lpr{M}(g)=\min_{m\in M}m\cdot g,\qquad
 \Dscr_M=\{g:\lpr{M}(g)>0\}\cup\{g\ge0:g\ne0\}.
\]
Also write $\lowprob{M}(\omega_i)=\min_{m\in M}m_i$ and
$\upprob{M}(\omega_i)=\max_{m\in M}m_i$.
Every coherent set agrees a.e. with some $\Dscr_M$; an elementary proof is in
Appendix~\ref{app:representation}.

Let $\mathcal C$ comprise all models coherent up to a null set. For a class
$\mathfrak R$ of rankings, let $\mathcal C_R$ be the coherent undominated
models evaluated by $R$. We retain the following requirements:
\begin{enumerate}[leftmargin=*,itemsep=3pt]
\item[\textbf{(E)}] Every model is evaluated by some $R\in\mathfrak R$.
\item[\textbf{(LA)}] $\mathcal C_R\ne\varnothing$ for every $R\in\mathfrak R$.
\item[\textbf{(A)}] Every member of $\mathcal C$ belongs to some $\mathcal C_R$.
\item[\textbf{(S)}] Each ranking admits a selection $p\mapsto C_p\in\mathcal C_R$
that is competitive on $\supp p$ and for which
$\{p:C_p\preceq_i\Dscr\}$ is closed for every $i,\Dscr$.
\item[\textbf{(C)}] A model tied at every world with a member of $\mathcal C_R$
also belongs to $\mathcal C_R$.
\end{enumerate}
All rankings constructed below evaluate every measurable model.

\section{Why the ordinal benchmark recipe fails independence}

The difficulty is structural. A model accepting every strictly negative
gamble cannot improve on any coherent benchmark at any world: those
acceptances are errors at every world. This holds for the original background
order and for \eqref{eq:advice-order}. In the ordinal recipe of v2, this leaves
no available index anywhere. Both comparison levels then take their default
value, and comparisons revert to the background order. Adding the same
contradictions to two otherwise consistent models can therefore erase an
earlier comparison.

The next proposition isolates this problem without depending on the
particular family, the maximum defining a score, or the size of a bonus.

\begin{proposition}[A common change can erase a comparison]\label{prop:obstruction}
Suppose a ranking on all measurable models has both properties below:
\begin{enumerate}[label=(\alph*)]
\item Every model not coherent modulo null sets is strictly dominated at
every world by a coherent model.
\item When both models accept all of $B=\R^n_{>0}\cup\R^n_{<0}$ and give
no contradictory advice outside $B$, their comparison at each world is exactly
$\unlhd_i$.
\end{enumerate}
Then the ranking fails \textbf{(IB)}.
\end{proposition}
\begin{proof}
Let $\Escr$ follow $D_1$ on $\{g:\|g\|_2<1\}$ and $D_2$ on its complement.
The radius condition is unchanged by $g\mapsto-g$, so $\Escr$ gives exactly
one recommendation on almost every pair. It follows the unanimously better
option when all payoffs have the same sign. It is not coherent even modulo
null sets: on directions with $g_1>0>g_2$, it accepts at radii below one and
rejects at radii above one. No cone can agree with this on almost every ray;
polar integration makes the disagreement positive in measure.

By (a), some coherent $C$ satisfies $C\prec_i\Escr$ at every world.
Nevertheless $C\unlhd_i\Escr$ cannot hold at every world. On any pair where
$C$ differs from $\Escr$, the latter gives a definite recommendation that is
good at at least one world (world 1 or world 2). At such a world, $C$ gives
silence or the opposite recommendation, so is worse. Since $C\not\approx\Escr$
and there are finitely many worlds, this happens on a positive-measure set
at some fixed world $i$.

Put $B=\R^n_{>0}\cup\R^n_{<0}$ and $A=B^c$. Both $C$ and $\Escr$ accept
exactly $\R^n_{>0}$ on $B$, up to null sets. Replace that common advice by
acceptance of all of $B$, giving $C'=C\cup B$ and $\Escr'=\Escr\cup B$.
Both now give contradictory advice on $B$ and nowhere else, up to null sets.
Their background comparison on
$A$ is unchanged, so $C'\not\unlhd_i\Escr'$ at the world just found. By (b),
$C'\not\preceq_i\Escr'$. Originally $C\prec_i\Escr$. This contradicts
\eqref{eq:independence}, with common outside advice changed from
$\R^n_{>0}$ to $B$.
\end{proof}

The original ordinal construction has both properties: its background order
agrees with $\unlhd_i$ outside $B$ on these otherwise consistent models, and
the models give identical advice on $B$. Updating it to the new background
order while retaining its strict dominance conclusion and default comparison
also leaves the obstruction intact.
This is not an impossibility theorem for independent rankings in general.
The next construction avoids the model-wide change of comparison rule.

\begin{remark}[Keeping the old indices does not give a local Bayes construction]
There is a second, narrower obstacle. For the three-world vacuous target,
the homothetic family in v2 has $K_p=\{p\}$ on the boundary and
$K_{(1/3,1/3,1/3)}=\Delta$. Take
\[
 p=(2/3,0,1/3),\qquad q=(0,2/3,1/3),\qquad g=(2,1,-1).
\]
Both $C_p$ and $C_q$ recommend $+$, whereas $C_{(p+q)/2}$ gives no advice.
These facts persist on an open neighbourhood of $g$. For fixed local
penalties, if $+$ uniquely minimizes expected penalty at both $p$ and $q$,
it also does so at their average: average its strict comparisons with each
other action. Thus this family cannot be the unique pointwise Bayes family
with those same probability indices. This does not rule out reindexing,
other families, or more general ordinal rankings.
\end{remark}

\section{Penalties for individual binary problems}

A direct way to ensure independence is to assign each piece of advice its own
penalty and add these penalties over problems. The penalty at a problem may
depend on its payoffs and the target, but never on the model's advice elsewhere.

For $g\in H$ and world $i$, choose measurable local penalties
$\ell_i(g,a)$, $a\in\{+,0,-,\times\}$, with
\begin{equation}
 \ell_i(g,\text{good})=0
 <\ell_i(g,0)
 <\ell_i(g,\text{bad})
 <\ell_i(g,\times)\le h_i(g),
 \qquad \int_Hh_i\,d\mu<\infty.
 \label{eq:local-penalties}
\end{equation}
These inequalities are required off the null coordinate hyperplanes.
Define
\begin{equation}
 L(\Dscr,\omega_i)=\int_H\ell_i(g,a_{\Dscr}(g))\,d\mu(g),
 \qquad
 \Dscr\preceq_i\Escr\ \Longleftrightarrow\
 L(\Dscr,\omega_i)\le L(\Escr,\omega_i).
 \label{eq:integral-loss}
\end{equation}
These losses are finite and uniformly bounded over models.

\begin{lemma}\label{lem:local-independence}
The ranking \eqref{eq:integral-loss} satisfies \textbf{(G)} and \textbf{(IB)}.
\end{lemma}
\begin{proof}
A guidance improvement decreases the integrand pointwise, strictly on a
positive-measure set when the improvement is strict. Its integral therefore
decreases accordingly. For independence, $A=-A$ means that a pair lies wholly
inside $A$ or wholly outside it. The common outside contribution cancels from
the difference of the two finite losses. Thus both comparisons in
\eqref{eq:independence} depend only on the same integral over $H\cap A$.
\end{proof}

It remains to make the undominated models coherent and to include the target.
The next theorem reduces both tasks to a pointwise calculation with four
numbers. A model is called $p$-Bayes if it minimizes
$\sum_i p_iL(\Dscr,\omega_i)$.

\begin{theorem}[A sufficient construction in any number of worlds]\label{thm:general}
Suppose \eqref{eq:local-penalties} holds. Suppose also that, for every
$p\in\Delta$, the minimizing advice in
\begin{equation}
 \min_{a\in\{+,0,-,\times\}}\sum_i p_i\ell_i(g,a)
 \label{eq:pointwise-bayes}
\end{equation}
is unique for almost every $g\in H$ and agrees a.e. with the advice of a
coherent model $C_p$. Then:
\begin{enumerate}[label=(\alph*)]
\item $C_p$ is the unique $p$-Bayes model modulo null sets and is undominated.
\item Each $C_p$ is competitive on $\supp p$; every map
$p\mapsto L(C_p,\omega_i)$ is continuous, giving the closedness in \textbf{(S)}.
\item Every model outside $\{C_p:p\in\Delta\}$ modulo null sets is strictly
dominated at every world by some $C_q$.
\end{enumerate}
The ranking satisfies \textbf{(LA)}, \textbf{(S)}, \textbf{(C)}, \textbf{(G)},
and \textbf{(IB)}, and evaluates every model. If such penalties can be chosen
through each coherent target, their class satisfies \textbf{(A)} as well.
\end{theorem}
\begin{proof}
For any $\Dscr$, subtract the pointwise minimum from its expected local loss
and integrate. The difference is nonnegative, and is positive unless the
advice agrees a.e. with that of $C_p$. Agreement of advice on almost every
pair is equivalent to agreement of acceptance sets a.e. Thus $C_p$ is uniquely
$p$-Bayes modulo null sets. A model dominating it would also be $p$-Bayes,
hence agree with it a.e. and tie everywhere, a contradiction.

For continuity, let $p_k\to p$. At almost every $g$, the unique minimizing
action at $p$ has a positive gap over the other three actions. Since their
expected penalties are linear in $p$, the minimizing action is eventually
unchanged. For each $k$ we may discard the null exceptional set in the
hypothesis; their countable union is still null. Dominated convergence, with
bound $h_i$, gives $L(C_{p_k},\omega_i)\to L(C_p,\omega_i)$.
Bayes optimality ensures that some $i\in\supp p$ has
$L(C_p,\omega_i)\le L(\Dscr,\omega_i)$, proving competitiveness.

For (c), fix $\Dscr$ outside the selected family modulo null sets. Uniqueness
gives a strictly positive, continuous function on the compact simplex:
\[
 d(p)=\sum_i p_i\bigl(L(\Dscr,\omega_i)-L(C_p,\omega_i)\bigr)>0.
\]
Let $\delta=\min_pd(p)>0$. The closed sets
\[
 A_i=\{p:L(C_p,\omega_i)\le L(\Dscr,\omega_i)-\delta/2\}
\]
cover each face in the required sense: for every $p$, some $i\in\supp p$
has $p\in A_i$, since the weighted average improvement is at least $\delta$.
The simplex lemma in Appendix~\ref{app:simplex} gives a common $q\in\bigcap_iA_i$.
Thus $C_q$ improves on $\Dscr$ at every world by at least $\delta/2$.

Finally, a model tied everywhere with $C_p$ is also $p$-Bayes and so agrees
with $C_p$ a.e., proving (C). The remaining requirements follow from the
properties just proved and Lemma~\ref{lem:local-independence}.
\end{proof}

This is a sufficient method, not a representation theorem for all rankings
satisfying independence. In particular, it places no restriction on what a
more general, incomplete ordinal ranking might achieve.

\section{Full coverage in two worlds}

Now let $n=2$, identifying a distribution with its first coordinate
$t\in[0,1]$. Every credal set is an interval $[l,u]\subseteq[0,1]$.
The whole class of rankings will have just one parameter, $w\in[0,1]$.
Larger $w$ makes silence cheaper. To cover a target interval, take $w$ to be
its width.

At a mixed-sign pair, write its representative in $H$ as $g=(a,-b)$ with
$a,b>0$. Give silence the same penalty at both worlds:
\begin{equation}
 \boxed{\quad c_w(a,b)=\frac{ab}{(1+w)(a+b)}.\quad}
 \label{eq:silence-costs}
\end{equation}
The complete table of penalty vectors is
\begin{equation}
 \begin{array}{c|cccc}
 \text{advice}&+&0&-&\times\\ \hline
 (\ell_1,\ell_2)&(0,b)&(c_w,c_w)&(a,0)&(2a,2b).
 \end{array}
 \label{eq:mixed-table}
\end{equation}
Thus good advice costs nothing; bad advice incurs the payoff shortfall;
contradictory advice twice that amount; and silence an intermediate penalty,
since $0<c_w<\min\{a,b\}$. At a unanimous pair $g=(a,b)$ with $a,b>0$, put
\begin{equation}
 \begin{array}{c|cccc}
 \text{advice}&+&0&-&\times\\ \hline
 (\ell_1,\ell_2)&(0,0)&(a/2,b/2)&(a,b)&(2a,2b).
 \end{array}
 \label{eq:unanimous-table}
\end{equation}
Values on the coordinate hyperplanes are immaterial.
Every penalty at world $i$ is at most $2|g_i|$, which is integrable under
Gaussian measure. Thus \eqref{eq:local-penalties} holds.

The next calculation explains the choice \eqref{eq:silence-costs}.
The two probabilities at which the best advice changes will become the
endpoints of a coherent credal interval.

\begin{lemma}[When each recommendation minimizes expected penalty]\label{lem:thresholds}
For $g=(a,-b)$ put $r=b/(a+b)$. Under prior $(t,1-t)$, the unique minimizing
advice, except at the displayed thresholds, is
\begin{equation}
 \begin{cases}
 -,&t<r/(1+w),\\
 0,&r/(1+w)<t<(r+w)/(1+w),\\
 +,&t>(r+w)/(1+w).
 \end{cases}\label{eq:thresholds}
\end{equation}
For $w=0$ the silence interval is empty and the thresholds coincide at $r$.
At a unanimous pair, $+$ uniquely minimizes expected penalty for every $t$.
\end{lemma}
\begin{proof}
The expected penalties for $-$, $0$, and $+$ are
\[
 ta,\qquad c_w(a,b),\qquad (1-t)b.
\]
The constant middle term meets the first at $t=r/(1+w)$ and the last at
$t=(r+w)/(1+w)$. The first and last terms meet at $t=r$, between these two
thresholds. This gives \eqref{eq:thresholds}, including its limiting case
$w=0$. Contradictory advice has a strictly greater penalty at both worlds
than any other choice. The unanimous case is immediate.
\end{proof}

The resulting family consists of intervals of width $w$, shortened where
necessary to remain in $[0,1]$:
\begin{equation}
 \boxed{\quad
 K_t=[l_w(t),u_w(t)]
 =\bigl[\max\{0,(1+w)t-w\},\ \min\{1,(1+w)t\}\bigr].
 \quad}\label{eq:interval-family}
\end{equation}
It is nonempty, since $l_w(t)\le t\le u_w(t)$.
Its lower and upper probabilities at $\omega_1$ are $l_w(t)$ and $u_w(t)$;
at $\omega_2$ they are $1-u_w(t)$ and $1-l_w(t)$.
Put $C_t=\Dscr_{K_t}$.

\begin{proposition}\label{prop:bayes-interval}
For every $t\in[0,1]$, the pointwise minimizing advice agrees a.e. with $C_t$
and is unique a.e. Consequently Theorem~\ref{thm:general} applies, including
at $t=0$, $t=1$, and $w=0$.
\end{proposition}
\begin{proof}
For $g=(a,-b)$,
\[
 \lpr{K_t}(g)=(a+b)l_w(t)-b,\qquad
 \lpr{K_t}(-g)=b-(a+b)u_w(t).
\]
Thus $C_t$ accepts $g$ exactly when $l_w(t)>r$, equivalent to
$t>(r+w)/(1+w)$. It accepts $-g$ exactly when $u_w(t)<r$, equivalent to
$t<r/(1+w)$. Between these thresholds it accepts neither. This agrees with
Lemma~\ref{lem:thresholds}. Every unanimous positive gamble is accepted.

For fixed $t$, a tie can occur only on a line with constant
$r=b/(a+b)$, or on the coordinate hyperplanes. These sets are Lebesgue-null,
also when $w=0$ and all three noncontradictory choices tie on the same line.
The thresholds lie strictly inside $(0,1)$, so the endpoint priors create no
additional ties. In particular, $C_0\approx D_2$ and $C_1\approx D_1$.
\end{proof}

The reference intervals have now been derived from the advice penalties.
The next theorem shows that every target appears in one of these families.

\begin{theorem}[Full two-world coverage with independence]\label{thm:full}
For two worlds, the class of finite loss rankings defined by
\eqref{eq:mixed-table}--\eqref{eq:unanimous-table}, as $w$ ranges over $[0,1]$,
satisfies \textbf{(E)}, \textbf{(LA)}, \textbf{(A)}, \textbf{(S)}, \textbf{(C)},
\textbf{(G)}, and \textbf{(IB)}. Every coherent set of desirable gambles is
undominated under some ranking in this class, including precise, boundary,
and vacuous targets.
\end{theorem}
\begin{proof}
Fix any target interval $M=[l,u]$, $0\le l\le u\le1$. Set
\begin{equation}
 \boxed{\quad w=u-l,\qquad t_0=\frac{u}{1+w}.\quad}
 \label{eq:target-parameters}
\end{equation}
Then $w\in[0,1]$, $t_0\in[0,1]$, and substitution in
\eqref{eq:interval-family} gives
\[
 l_w(t_0)=u-w=l,\qquad u_w(t_0)=u.
\]
Hence $C_{t_0}=\Dscr_M$. Proposition~\ref{prop:bayes-interval} and
Theorem~\ref{thm:general} give all the ranking-specific requirements, and
make the target undominated. This includes singletons ($w=0$), intervals
containing either vertex, and the vacuous interval ($w=1$).

Every coherent set agrees a.e. with $\Dscr_M$ for some interval $M$ by
Appendix~\ref{app:representation}. Thus the class satisfies (A). Every ranking
evaluates every measurable model, giving (E).
\end{proof}

\paragraph{Why the target is undominated.}
At the prior $t_0$ in \eqref{eq:target-parameters}, its recommendation minimizes
expected penalty in almost every binary problem. Any different recommendation
on a positive-measure set raises the expected total penalty. The pointwise
uniqueness also rules out domination when the target prior is an endpoint.

\paragraph{The vacuous target.}
For $M=[0,1]$, take $w=1$ and $t_0=1/2$. Then
\begin{equation}
 c_1(a,b)=\frac{ab}{2(a+b)},\qquad
 K_t=[\max\{0,2t-1\},\ \min\{1,2t\}].
 \label{eq:vacuous}
\end{equation}
At the uniform prior, the penalty for silence is strictly below both $a/2$
and $b/2$. So silence uniquely minimizes expected penalty in every
mixed-sign problem. In unanimous problems, choosing the better option
uniquely minimizes it. These are exactly the recommendations of the vacuous
coherent set, modulo null sets.

\section{Three or more worlds: an obstruction to additive penalties}

The two-world result does not extend to this whole additive method. Allowing
silence penalties to depend on the world, or changing the penalties for bad
advice, cannot restore full coverage. The obstruction also allows ties: it
does not assume the uniqueness hypothesis of Theorem~\ref{thm:general}.

\begin{theorem}[No full additive coverage in three or more worlds]\label{thm:additive-obstruction}
Let $n\ge3$. A ranking of all measurable models with losses
\eqref{eq:integral-loss} and strictly ordered local penalties
\eqref{eq:local-penalties} cannot satisfy \textbf{(LA)}, \textbf{(S)}, and
make the vacuous model $\Dscr_\Delta$ undominated.
In fact, if \textbf{(LA)} and \textbf{(S)} hold, every $\Dscr_M$ whose credal
set contains three distinct simplex vertices is strictly dominated at every
world by a coherent model.
Consequently no class of these rankings satisfies full \textbf{(A)}.
\end{theorem}

The proof, given in Appendix~\ref{app:obstruction}, proceeds as follows.
\begin{itemize}[leftmargin=*,itemsep=2pt]
\item The axioms require a coherent minimum of expected loss at every prior.
\item An undominated target must minimize expected loss at some prior.
\item A target containing three vertices forces silence on each associated
two-world face, at suitable priors.
\item Those faces impose incompatible requirements on a coherent model when
a third payoff changes sign.
\end{itemize}

The last step explains the dimensional difference. With only two worlds,
there is no third payoff across which the separate trade-offs must fit
together. The interval construction exploits exactly that freedom.

This result covers arbitrary measurable, world-dependent local penalties,
not just payoff-shortfall penalties or a common cost of silence. It assumes
neither continuity nor homogeneity of those penalties and does not prescribe
a credal family or its indexing. It also does not use \textbf{(C)}.

\paragraph{Allowing infinite additive losses does not avoid this result.}
For nonnegative local penalties on all measurable models, strict
\textbf{(G)} forces the penalty of the everywhere-contradictory model to be
finite at each world. Otherwise split its infinite penalty between two
measurable sets of pairs, each still having infinite penalty. Correcting its
advice on one set is a strict guidance improvement, but leaves the total
penalty infinite. This violates \textbf{(G)}. Such a split exists because
the measure on binary problems has no atoms; when the density is finite,
one can partition it into finite-mass pieces and distribute them into two
divergent sums, and an infinite-valued part of positive measure can itself
be split in two. The worst local penalty is therefore an integrable envelope,
as already required in \eqref{eq:local-penalties}.

\paragraph{What remains open.}
Independence \textbf{(IB)} alone has not been shown to imply an additive
representation. The theorem therefore leaves open full coverage by more
general independent ordinal rankings, and by nonadditive real-valued or
extended-real-valued losses. It does not show that incomparabilities are
necessary. What it does establish is that changing the local penalty formula
cannot solve the general problem within the additive approach.

There are still positive results in every dimension. Penalties
$0,|g_i|/2,|g_i|,2|g_i|$ for good, no, bad, and contradictory advice yield
the precise Bayes family $C_p\approx\{g:p\cdot g>0\}$, so
Theorem~\ref{thm:general} applies. The remaining coverage of imprecise targets
within the additive class is not classified here.

\appendix
\section{Representing coherent sets by credal sets}\label{app:representation}

For completeness, this connects coverage of credal targets to coverage of all
coherent sets, including coherent sets with different boundary conventions.

\begin{lemma}\label{lem:representation}
For every coherent $\Dscr\subseteq\R^n$, the set
\[
 M=\{m\in\Delta:m\cdot g\ge0\text{ for every }g\in\Dscr\}
\]
is a credal set, and
$\{\lpr{M}>0\}\subseteq\Dscr\subseteq\{\lpr{M}\ge0\}$.
Consequently $\Dscr$ is Lebesgue measurable and $\Dscr\approx\Dscr_M$.
Conversely $\Dscr_M$ is coherent for every credal set $M$.
\end{lemma}
\begin{proof}
Let $C=\overline\Dscr$. It is a closed convex cone containing zero.
No strictly negative vector lies in $C$: coherence excludes every nonzero
$g\le0$, since adding $-g\in\Dscr$ would put zero in $\Dscr$.
A strictly negative vector has a neighbourhood of excluded vectors.

If $h\notin C$, choose $c\in C$ nearest to $h$ in Euclidean distance and put
$w=c-h\ne0$. Minimality along segments in $C$ gives
$w\cdot(x-c)\ge0$ for all $x\in C$ (expand the squared distance along the
segment and let its length tend to zero). Taking $x=0$ and $x=2c$ gives
$w\cdot c=0$. Hence $w\cdot x\ge0$ on $C$ and
$w\cdot h=-\|w\|_2^2<0$. Since $e_i\in\Dscr$, each $w_i\ge0$.
Normalization gives $m=w/\sum_iw_i\in M$ with $m\cdot h<0$.
Applying this to $h=-\one$ proves that $M$ is nonempty. It is closed and convex
by definition.

The inclusion $\Dscr\subseteq\{\lpr{M}\ge0\}$ is immediate.
If $\lpr{M}(g)=2\varepsilon>0$, put $h=g-\varepsilon\one$. Since
$\lpr{M}(h)=\varepsilon>0$, the separation just proved gives $h\in C$.
Choose $h'\in\Dscr$ with $\|h'-h\|_\infty<\varepsilon$. Every coordinate of
$g-h'$ is positive, so coherence gives $g=h'+(g-h')\in\Dscr$.

The lower prevision is continuous, since
$|\lpr{M}(g)-\lpr{M}(h)|\le\|g-h\|_\infty$, and
$\lpr{M}(g+t\one)=\lpr{M}(g)+t$. Its zero set therefore meets each line
parallel to $\one$ once and is null by Fubini's theorem. Both $\Dscr$ and
$\Dscr_M$ differ from the open set $\{\lpr{M}>0\}$ only within this null set,
proving measurability and agreement a.e.

Conversely, positive homogeneity and
$\lpr{M}(g+h)\ge\lpr{M}(g)+\lpr{M}(h)$ show that $\Dscr_M$ is closed under
positive scaling and addition; it includes the nonzero nonnegative gambles
and excludes zero by definition. Thus it is coherent.
\end{proof}

\section{The simplex lemma used for strict dominance}\label{app:simplex}

\begin{lemma}\label{lem:kkm}
Suppose closed sets $A_1,\ldots,A_n\subseteq\Delta$ have the property that,
for every $p\in\Delta$, some $i\in\supp p$ satisfies $p\in A_i$.
Then $\bigcap_iA_i\ne\varnothing$.
\end{lemma}
\begin{proof}
Each $A_i$ contains $e_i$ and is nonempty. Let $d_i(p)$ be the Euclidean
distance from $p$ to $A_i$ and let $s(p)=\sum_i d_i(p)$.
The map $F:\Delta\to\Delta$ defined by
\[
 F_i(p)=\frac{p_i+d_i(p)}{1+s(p)}
\]
is continuous. Brouwer's fixed-point theorem supplies $F(p)=p$.
Thus $p_i s(p)=d_i(p)$ for every $i$. By the hypothesis some $i$ has
$p_i>0$ and $d_i(p)=0$, so $s(p)=0$. Every $d_i(p)$ is therefore zero,
and closedness gives $p\in\bigcap_iA_i$.
\end{proof}

For two worlds, only the elementary interval case is needed: two closed sets
covering $[0,1]$, one containing $0$ and the other containing $1$, must intersect.
Otherwise they would partition the interval into disjoint nonempty closed
sets, contradicting its connectedness.

\section{Why additive full coverage fails for three or more worlds}\label{app:obstruction}

There are two parts to the proof. First, the axioms turn the coverage question
into a question about expected-loss minima. Then a three-world argument
shows that the required minima cannot all be coherent.

\subsection*{From the axioms to coherent expected-loss minima}

\begin{lemma}\label{lem:bayes-necessary}
For losses \eqref{eq:integral-loss} satisfying
\eqref{eq:local-penalties}:
\begin{enumerate}[label=(\alph*)]
\item If \textbf{(LA)} and \textbf{(S)} hold, a coherent $p$-Bayes model
exists for every $p\in\Delta$.
\item If a model is not strictly dominated at every world, it is $p$-Bayes
for some $p\in\Delta$.
\item If $\Dscr_M$ is $p$-Bayes, then $M\subseteq\Delta_{\supp p}$, where
$\Delta_S=\{m\in\Delta:m_j=0\text{ for }j\notin S\}$.
\end{enumerate}
\end{lemma}
\begin{proof}
(a) Choose a measurable pointwise minimizing action in
\eqref{eq:pointwise-bayes}, breaking ties in a fixed order, and call the
resulting model $B$. It is $p$-Bayes. Apply Lemma~\ref{lem:kkm} to
\[
 A_i=\{q:C_q\preceq_i B\},
\]
using the selection in (S). Closedness and the covering condition are
exactly the two parts of (S). Some coherent $C_q$ therefore does at least
as well as $B$ at every world. It must also be $p$-Bayes.

(b) Let $\mathcal L\subseteq\R^n$ be the attainable loss vectors. Its closure
is convex. Here is the needed elementary approximation argument. Given
two advice rules and $0\le t\le1$, approximate their vector of local loss
differences in $L^1$ by a simple function. In each cell of its finite
partition, choose a measurable subset with proportion $t$ of the cell's
measure. Such subsets exist because Gaussian measure has no atoms. Follow
the first rule on these subsets and the second elsewhere. Its loss vector
approaches the weighted average of the two original loss vectors as the
approximation error tends to zero. The same conclusion for limits proves
convexity of $\overline{\mathcal L}$.

Let $v$ be the loss vector of the model in (b).
The convex sets $\overline{\mathcal L}$ and $v-(0,\infty)^n$ are disjoint:
a point of their intersection could be approximated by an attainable vector
strictly below $v$ in every coordinate. The separating-hyperplane theorem
therefore gives a nonzero $p\ge0$ with
\[
 p\cdot x\ge p\cdot v\qquad(x\in\overline{\mathcal L}).
\]
The coefficients are nonnegative because $v-(0,\infty)^n$ is unbounded in
every negative coordinate direction. Normalize them to sum to one. This
is exactly $p$-Bayes optimality.

(c) Write $S=\supp p$. Whenever $g_i>0$ for every $i\in S$, accepting $g$
is uniquely best in expected local penalty. A coherent $p$-Bayes model must
therefore accept almost every such gamble, whatever its other coordinates.
If some $m\in M$ had $m_j>0$ outside $S$, setting the coordinates in $S$
positive and $g_j$ sufficiently negative would give an open set of these
gambles with $m\cdot g<0$. None can be accepted by $\Dscr_M$, a contradiction.
Thus $M\subseteq\Delta_S$. In particular, the model's advice depends only
on the coordinates in $S$, up to null sets.
\end{proof}

\subsection*{The compatibility requirement imposed by a third world}

The following lemma isolates the geometric difficulty. A positive contrast
means that accepting a gamble beats silence. The lemma does not assume
that the contrasts are continuous or homogeneous. Coherence itself forces
the two-world trade-offs to depend only on payoff ratios.

\begin{lemma}[Three-world compatibility]\label{lem:three-world}
Let $I$ contain at least three world indices, and write
$p\cdot\psi=\sum_{i\in I}p_i\psi_i$. Let measurable functions
$\psi_i(g)$, $i\in I$, be finite and satisfy
$\operatorname{sgn}\psi_i(g)=\operatorname{sgn}g_i$ off null sets.
Suppose that for every $p\in\Delta_I$ there is a coherent $\Dscr_{K_p}$ with
$K_p\subseteq\Delta_{\supp p}$ such that, up to null sets,
\begin{equation}
 \{g:p\cdot\psi(g)>0\}\ \subseteq\ \Dscr_{K_p}
 \subseteq\ \{g:p\cdot\psi(g)\ge0\}.
 \label{eq:contrast-completion}
\end{equation}
For each distinct $i,j\in I$, on $\{g_i>0>g_j\}$ there is a positive,
finite, nondecreasing function $F_{ij}:(0,\infty)\to(0,\infty)$ such that
\begin{equation}
 \frac{\psi_i(g)}{-\psi_j(g)}
 =F_{ij}\!\left(\frac{g_i}{-g_j}\right)\quad\text{a.e.},
 \qquad \lim_{t\to\infty}F_{ij}(t)=\infty.
 \label{eq:edge-ratio}
\end{equation}
In particular, that contrast ratio cannot have a finite upper bound.
\end{lemma}
\begin{proof}
\emph{First obtain the functions $F_{ij}$.}
Fix $i\ne j$ and restrict to $g_i>0>g_j$. Put
$R(g)=\psi_i(g)/(-\psi_j(g))$.
For weights proportional to $(1,s)$ on these two worlds,
\eqref{eq:contrast-completion} says
\[
 \{R>s\}\subseteq\Dscr_{K_s}\subseteq\{R\ge s\}
 \quad\text{on this region, a.e.}
\]
Since $K_s$ lies on the two-world edge, its acceptance set here is, up to a
null set, an upper interval in the single ratio $g_i/(-g_j)$; it is
independent of all other payoffs. Ties cause no difficulty: for each
rational $t>0$,
\[
 \{R>t\}=\bigcup_{\substack{s\in\mathbb Q\\s>t}}\Dscr_{K_s}
 \quad\text{on this region, a.e.}
\]
Every set on the right is an upper interval in that payoff ratio.
The rational strict level sets therefore show that
$R(g)=F_{ij}(g_i/(-g_j))$ a.e. for a nondecreasing $F_{ij}$.
It is positive and finite: a zero or infinite value at an interior argument
would, by monotonicity, make $R$ zero or infinite on a positive-measure set.

\emph{Now rule out a finite upper bound.}
Suppose $d=\lim_{t\to\infty}F_{ij}(t)<\infty$, and choose $k\in I$ distinct
from $i,j$. Initially suppose $F_{ik}$ is nonconstant. Write
$B=\lim_{t\to\infty}F_{ik}(t)$, allowing $B=\infty$.
Choose a continuity point $t>0$ of this monotone function with
$a=F_{ik}(t)<B$. For a sufficiently small $\varepsilon>0$, the positive
weights
\begin{equation}
 q_i=1,\qquad q_j=\varepsilon,\qquad
 q_k=a\left(1-\frac{\varepsilon}{2d}\right)
 \quad\text{satisfy}\quad
 1>\frac{\varepsilon}{d}+\frac{q_k}{B},
 \label{eq:three-weights}
\end{equation}
where $q_k/B=0$ if $B=\infty$. Apply the hypothesis at the prior
$p=q/(q_i+q_j+q_k)$, with zero weight elsewhere. We compute signs using
the unnormalised weights $q$, which gives the same signs as $p$.

Near the three-payoff vector $(g_i,g_j,g_k)=(t,0,-1)$, the sign of
$q\cdot\psi$ changes incompatibly across $g_j=0$. On the positive side,
divide by $-\psi_k$ to obtain
\[
 \frac{q\cdot\psi}{-\psi_k}
 =F_{ik}\!\left(\frac{g_i}{-g_k}\right)
   +\varepsilon F_{jk}\!\left(\frac{g_j}{-g_k}\right)-q_k>0.
\]
This holds throughout a sufficiently small positive-side neighbourhood,
since $a-q_k=\varepsilon a/(2d)>0$.
On the negative side, divide by $\psi_i$:
\[
 \frac{q\cdot\psi}{\psi_i}
 =1-\frac{\varepsilon}{F_{ij}(g_i/(-g_j))}
     -\frac{q_k}{F_{ik}(g_i/(-g_k))}<0.
\]
Indeed $F_{ij}\le d$, and at $g_i/(-g_k)=t$ the right side is at most
$1-\varepsilon/d-q_k/a=-\varepsilon/(2d)$. Continuity at $t$ preserves
the strict inequality nearby. These statements hold a.e.; all additional
world coordinates are immaterial to them.

Yet for $(g_i,g_j,g_k)=(T,-1,-1)$ with $T$ sufficiently large,
\eqref{eq:three-weights} gives
\[
 1-\frac{\varepsilon}{F_{ij}(T)}-\frac{q_k}{F_{ik}(T)}>0.
\]
The inequality also holds on a neighbourhood. A coherent cone cannot
contain the positive-side neighbourhood near $(t,0,-1)$ and this additional
point with $g_j<0$, while excluding the adjacent negative-side neighbourhood.
Taking a small convex combination with the additional point crosses into
that excluded neighbourhood. The contradiction is unchanged by null sets:
the interior of a coherent cone contains every open set it contains a.e.,
and cannot meet an open set it excludes a.e.

If $F_{ik}$ is constant but $F_{ij}$ is not, interchange $j$ and $k$ in the
argument. Both relevant upper bounds are now finite, so it applies again.
Finally, if $F_{ij}=A$ and $F_{ik}=B$ are both constant, use weights
\[
 q_i=1,\qquad q_j=3A/4,\qquad q_k=3B/4.
\]
The contrast is strictly positive whenever $g_i,g_j>0>g_k$, and whenever
$g_i,g_k>0>g_j$, but strictly negative whenever $g_i>0>g_j,g_k$.
For example, in the three indicated coordinates, it accepts
$(1,2,-3)$ and $(1,-3,2)$ but excludes their sum $(2,-1,-1)$, with the same
strict signs on neighbourhoods. This again contradicts coherence.
All cases rule out $d<\infty$.
\end{proof}

\subsection*{Why a target containing three vertices violates that requirement}

We now apply the lemma to arbitrary local advice penalties. Extend the
penalties from $H$ to both orientations of each pair. Write $c_i(g)$ for
the penalty of silence and $b_i(g)$ for that of bad advice. Thus
$0<c_i(g)<b_i(g)$ and both functions are unchanged by $g\mapsto-g$.
At prior $p$, the three relevant expected penalties are
\begin{equation}
 V=\sum_{g_i<0}p_i b_i\quad\text{for }+,\qquad
 C=\sum_i p_i c_i\quad\text{for }0,\qquad
 U=\sum_{g_i>0}p_i b_i\quad\text{for }-.
 \label{eq:three-costs}
\end{equation}
Contradictory advice is always strictly worse. Minimizing total expected
loss is equivalent to minimizing these local penalties a.e.: integrate
the nonnegative difference from their pointwise minimum.

\begin{proposition}\label{prop:bayes-vertex-obstruction}
Suppose a coherent $p$-Bayes model exists for every $p\in\Delta$.
No coherent $\Dscr_M$ with three distinct simplex vertices in $M$ can be
$p^0$-Bayes at any prior $p^0$.
\end{proposition}
\begin{proof}
Suppose otherwise, and let $I$ be three indices whose vertices lie in $M$.
Lemma~\ref{lem:bayes-necessary}(c) gives $p_i^0>0$ for $i\in I$.
For each distinct $i,j\in I$, consider the two-world prior
\[
 r^{ij}=\frac{p_i^0e_i+p_j^0e_j}{p_i^0+p_j^0}.
\]
Any coherent $r^{ij}$-Bayes model is the vacuous model on this edge.
To see this, first take gambles with $g_j<0$ and every other coordinate
positive. The target gives silence, since $e_j\in M$ prevents acceptance
of $g$ and $e_i\in M$ prevents acceptance of $-g$.
Its $p^0$-Bayes optimality gives
\[
 \sum_{h\ne j}p_h^0c_h\le p_j^0(b_j-c_j).
\]
There is a third index in $I$ with positive weight and positive silence
penalty, so
$p_i^0c_i<p_j^0(b_j-c_j)$.
At $r^{ij}$, accepting $g$ is therefore strictly worse than silence on
this entire sign region, a.e.
By Lemma~\ref{lem:bayes-necessary}(c), a coherent $r^{ij}$-Bayes model
depends only on $(g_i,g_j)$. It consequently accepts none of the region
$g_i>0>g_j$, up to null sets, whatever the other payoffs.
Interchanging $i,j$ excludes acceptance on $g_j>0>g_i$ as well.
On unanimous two-world payoffs the better recommendation is uniquely best.
Thus its credal set is exactly the edge $\Delta_{\{i,j\}}$.

In particular, silence is $r^{ij}$-Bayes on every problem with
$g_i>0>g_j$, up to null sets. Comparing it with both recommendations gives
\begin{equation}
 \frac{c_i}{b_j-c_j}\le\frac{p_j^0}{p_i^0},
 \qquad
 \frac{c_i}{b_i}+\frac{c_j}{b_j}\le1.
 \label{eq:forced-edge-bounds}
\end{equation}
For the second inequality, multiply
$p_i^0c_i\le p_j^0(b_j-c_j)$ and
$p_j^0c_j\le p_i^0(b_i-c_i)$, then cancel the positive weights.
These conclusions apply to every opposite-sign pair of indices in $I$.

Define the acceptance-versus-silence contrasts
\begin{equation}
 \psi_i(g)=
 \begin{cases}
 c_i(g),&g_i>0,\\
 -(b_i(g)-c_i(g)),&g_i<0,
 \end{cases}
 \qquad i\in I.
 \label{eq:silence-contrast}
\end{equation}
They are strictly sign-preserving. For any $p\in\Delta_I$,
$p\cdot\psi=C-V$. If this is negative, accepting $g$ is not Bayes.
If it is positive, accepting $g$ uniquely minimizes expected penalty.
The latter assertion needs a check against the other recommendation.
When the supported payoffs have both signs, the second inequality in
\eqref{eq:forced-edge-bounds} allows a number $0<\lambda<1$ with
\[
 c_i\le\lambda b_i\quad(g_i>0, i\in\supp p),\qquad
 c_j\le(1-\lambda)b_j\quad(g_j<0, j\in\supp p).
\]
For example choose $\lambda$ between the largest positive-side $c_i/b_i$
and one minus the largest negative-side $c_j/b_j$.
Hence $C\le\lambda U+(1-\lambda)V$. If $V<C$, then also $V<U$.
For unanimous supported payoffs the assertion is immediate.

Every coherent $p$-Bayes model, $p\in\Delta_I$, therefore satisfies
\eqref{eq:contrast-completion}; its credal set lies in
$\Delta_{\supp p}$. Lemma~\ref{lem:three-world} applies.
But the first inequality in \eqref{eq:forced-edge-bounds} bounds
$\psi_i/(-\psi_j)$ by $p_j^0/p_i^0$ on $g_i>0>g_j$.
This is the contradiction.
\end{proof}

\begin{proof}[Proof of Theorem~\ref{thm:additive-obstruction}]
By Lemma~\ref{lem:bayes-necessary}(a), (LA) and (S) supply coherent Bayes
models at all priors. Proposition~\ref{prop:bayes-vertex-obstruction} rules
out Bayes optimality of the stated targets at every prior.
Lemma~\ref{lem:bayes-necessary}(b) therefore gives a model strictly better
than the target at every world. The same simplex argument used in (a)
gives a coherent model weakly better than this model at every world,
hence strictly better than the target everywhere.
Since $\Delta$ contains three vertices when $n\ge3$, the vacuous case follows.
\end{proof}

\end{document}
